% GENERATED FILE. Edit canonical lessons, then run npm run generate:books.
% canonical-source-hash: fnv1a64:bd941af4f3f6a0b9
% canonical-lessons: ZH-C369-shuoming, ZH-C369-danshi, ZH-W369-dan-but, ZH-C369-juli, ZH-C369-qici, ZH-C369-lingwai

\chapter{Making It Clear}
\label{ch:zh-making-it-clear}
\hlchaptermodality{369}

\begin{chapteropening}
\textbf{By the end of this chapter:} \emph{I can make something clear, give an example, and add a second point and a further one.}

Complete the last lesson of chapter 369: I can make something clear, give an example, and add a second point and a further one.
\end{chapteropening}

\section[\texorpdfstring{\zh{说明} (shuōmíng)}{说明 (shuōmíng)}]{\zh{说明} (shuōmíng) — to make something clear; to explain clearly}

\label{lesson:ZH-C369-shuoming}



\begin{quote}
Before the new one: say the Mandarin for to encourage, then the Mandarin for to make a choice.
\end{quote}

\subsection*{You'll want to know: \zh{说明}}
\textbf{\zh{说明}} — \emph{shuōmíng} — \textquotedblleft{}to make something clear; to explain clearly\textquotedblright{}.

\begin{quote}
\textbf{\zh{你必须说明为什么不能准时来}} \emph{nǐ bìxū shuōmíng wèishénme bù néng zhǔnshí lái} — You have to explain why you couldn't be on time.
\end{quote}

\begin{cousinweb}
\textbf{\zh{说明}} (\emph{shuōmíng}) is built from \zh{说} 'to say' and \zh{明} 'bright', both inherited from Middle Chinese, \emph{sywet} and \emph{mjaeng} (Guangyun, in Baxter's transcription), which glosses the first 'to tell' and the second 'bright'.
\end{cousinweb}

\begin{grammarlens}[title={What you've built}]
One more word for giving reasons.
\end{grammarlens}

\subsection*{Your turn}
\begin{itemize}
\raggedright
  \item \emph{Say it:} \emph{shuōmíng}
  \item \emph{Say it:} \emph{shuōmíng}, once more
  \item \emph{Say it:} \emph{nǐ bìxū shuōmíng wèishénme bù néng zhǔnshí lái} — You have to explain why you couldn't be on time.
  \item \emph{Recall:} say the Mandarin for to encourage, then the Mandarin for to make a choice, then say \emph{shuōmíng} again
\end{itemize}

\begin{tcolorbox}[breakable,colback=teal!4,colframe=teal!35!black,title={Before you move on}]
What does \zh{说明} mean? (\textquotedblleft{}To make something clear; to explain clearly\textquotedblright{}.) Say it once more.
\end{tcolorbox}
\section[\texorpdfstring{\zh{但是} (dànshì)}{但是 (dànshì)}]{\zh{但是} (dànshì) — but; however}

\label{lesson:ZH-C369-danshi}



\begin{quote}
Before the new one: say the Mandarin for to make a choice, then the Mandarin for to make something clear.
\end{quote}

\subsection*{You'll want to know: \zh{但是}}
\textbf{\zh{但是}} — \emph{dànshì} — \textquotedblleft{}but; however\textquotedblright{}.

\begin{quote}
\textbf{\zh{我觉得他很好但是我不喜欢他}} \emph{wǒ juéde tā hěn hǎo dànshì wǒ bù xǐhuan tā} — I think he's nice, but I don't like him.
\end{quote}

The new character, in pieces: \zh{但} — seven strokes in two pieces: the person piece on the left, two strokes; on the right \zh{日} \emph{rì}, four strokes, over a long bar, five strokes in all. The next lesson puts it in your hand.

\begin{cousinweb}
\textbf{\zh{但是}} (\emph{dànshì}) is built from \emph{dàn} and \zh{是}, both inherited from Middle Chinese, \emph{danX} and \emph{dzyeX} (Guangyun, in Baxter's transcription), which glosses the first as a particle and 'only', and the second 'right, as against wrong'. How the pair came to mean but is not settled.
\end{cousinweb}

\begin{grammarlens}[title={What you've built}]
One more word for giving reasons.
\end{grammarlens}

\subsection*{Your turn}
\begin{itemize}
\raggedright
  \item \emph{Say it:} \emph{dànshì}
  \item \emph{Say it:} \emph{dànshì}, once more
  \item \emph{Say it:} \emph{wǒ juéde tā hěn hǎo dànshì wǒ bù xǐhuan tā} — I think he's nice, but I don't like him.
  \item \emph{Recall:} say the Mandarin for to make a choice, then the Mandarin for to make something clear, then say \emph{dànshì} again
\end{itemize}

\begin{tcolorbox}[breakable,colback=teal!4,colframe=teal!35!black,title={Before you move on}]
What does \zh{但是} mean? (\textquotedblleft{}But; however\textquotedblright{}.) Say it once more.
\end{tcolorbox}
\section[\texorpdfstring{\zh{但} (dàn)}{但 (dàn)}]{\zh{但} — the person piece, then the sun over a line}

\label{lesson:ZH-W369-dan-but}



\begin{quote}
The word first, then the character inside it.

\begin{itemize}
\raggedright
  \item \emph{Recall:} say \emph{dànshì} — \textbf{R1}, one lesson back, before it decays
\end{itemize}

\begin{quote}
\zh{你} \zh{日}
\end{quote}

\emph{Read it:} each one aloud, with nothing said first — \emph{nǐ}, \emph{rì}
\end{quote}

\subsection*{Script — the person piece, then the sun over a line}
\hlblockfigure{\detokenize{figures/ZH-W369-dan-but-filmstrip.pdf}}{How \zh{但} is written, stroke by stroke}

\begin{quote}
\zh{但}
\end{quote}

\textbf{Two pieces, left before right.} Count the pieces, not the strokes.

\textbf{The left piece, \zh{亻}, two strokes}, as in \zh{你}: a left-falling stroke, then a vertical.

\textbf{The right piece, five strokes}: \zh{日}, four strokes, a vertical, a horizontal that turns down, a short horizontal inside and the bottom one; then \zh{一}, a long horizontal under it, with a little space between.

2 + 5 = seven.

\subsection*{Your turn}
\begin{itemize}
\raggedright
  \item \emph{Write it:} \zh{但}, all seven, with the model in front of you
  \item \emph{Cover:} write it again from memory, and count the pieces, not the strokes
\end{itemize}

\begin{tcolorbox}[breakable,colback=teal!4,colframe=teal!35!black,title={Before you move on}]
How many strokes in \zh{但}? (\textbf{Seven}.) What sits on the right? (\textbf{\zh{日} over \zh{一}}.) Which piece comes first? (\textbf{\zh{亻}, the person piece}.)

Source: \href{https://raw.githubusercontent.com/chanind/hanzi-writer-data/68d10a4b21150cae5e1ebbd223eed289cf32d90c/data/\%E4\%BD\%86.json}{Hanzi Writer Data \zh{但}, pinned stroke paths and medians}.
\end{tcolorbox}
\section[jǔlì]{jǔlì — to give an example}

\label{lesson:ZH-C369-juli}



\begin{quote}
Before the new one: say the Mandarin for to make something clear, then the Mandarin for however.
\end{quote}

\subsection*{You'll want to know: jǔlì}
\emph{jǔlì} — \textquotedblleft{}to give an example\textquotedblright{}.

\begin{quote}
\emph{nǐ néng jǔlì shuōmíng ma} — Can you give an example?
\end{quote}

Its characters come later; for now it is a word for the ear.

\begin{cousinweb}
\emph{jǔlì} is built from \emph{jǔ} and \emph{lì}, both inherited from Middle Chinese, \emph{kjoX} and \emph{ljejH} (Guangyun, in Baxter's transcription), which glosses the first 'to lift' and the second 'to compare'.
\end{cousinweb}

\begin{grammarlens}[title={What you've built}]
One more word for giving reasons.
\end{grammarlens}

\subsection*{Your turn}
\begin{itemize}
\raggedright
  \item \emph{Say it:} \emph{jǔlì}
  \item \emph{Say it:} \emph{jǔlì}, once more
  \item \emph{Say it:} \emph{nǐ néng jǔlì shuōmíng ma} — Can you give an example?
  \item \emph{Recall:} say the Mandarin for to make something clear, then the Mandarin for however, then say \emph{jǔlì} again
\end{itemize}

\begin{tcolorbox}[breakable,colback=teal!4,colframe=teal!35!black,title={Before you move on}]
What does jǔlì mean? (\textquotedblleft{}To give an example\textquotedblright{}.) Say it once more.
\end{tcolorbox}
\section[qícì]{qícì — secondly; second in order}

\label{lesson:ZH-C369-qici}



\begin{quote}
Before the new one: say the Mandarin for however, then the Mandarin for to give an example.
\end{quote}

\subsection*{You'll want to know: qícì}
\emph{qícì} — \textquotedblleft{}secondly; second in order\textquotedblright{}.

\begin{quote}
\emph{shǒuxiān yīnggāi zhǔnshí qícì yīnggāi yǒu lǐmào} — First, be on time; second, be polite.
\end{quote}

Its characters come later; for now it is a word for the ear.

\begin{cousinweb}
\emph{qícì} is built from \emph{qí} and \emph{cì}, both inherited from Middle Chinese, \emph{gi} and \emph{tshijH} (Guangyun, in Baxter's transcription), which glosses the first as a particle and the second 'order' and 'the next'.
\end{cousinweb}

\begin{grammarlens}[title={What you've built}]
One more word for giving reasons.
\end{grammarlens}

\subsection*{Your turn}
\begin{itemize}
\raggedright
  \item \emph{Say it:} \emph{qícì}
  \item \emph{Say it:} \emph{qícì}, once more
  \item \emph{Say it:} \emph{shǒuxiān yīnggāi zhǔnshí qícì yīnggāi yǒu lǐmào} — First, be on time; second, be polite.
  \item \emph{Recall:} say the Mandarin for however, then the Mandarin for to give an example, then say \emph{qícì} again
\end{itemize}

\begin{tcolorbox}[breakable,colback=teal!4,colframe=teal!35!black,title={Before you move on}]
What does qícì mean? (\textquotedblleft{}Secondly; second in order\textquotedblright{}.) Say it once more.
\end{tcolorbox}
\section[lìngwài]{lìngwài — besides; in addition}

\label{lesson:ZH-C369-lingwai}



\begin{quote}
Before the new one: say the Mandarin for to give an example, then the Mandarin for secondly.
\end{quote}

\subsection*{You'll want to know: lìngwài}
\emph{lìngwài} — \textquotedblleft{}besides; in addition\textquotedblright{}.

\begin{quote}
\emph{zhège tài guì le lìngwài wǒ yě bù xǐhuan} — This one's too expensive, and besides, I don't like it.
\end{quote}

Its characters come later; for now it is a word for the ear.

\begin{cousinweb}
\emph{lìngwài} is built from \emph{lìng} 'separate' and \zh{外} 'outside'. The second is inherited from Middle Chinese \emph{ngwajH} (Guangyun, in Baxter's transcription), glossed 'the outside' and 'far'. No dictionary we use records an older form of the first, so its history is not settled.
\end{cousinweb}

\begin{grammarlens}[title={What you've built}]
One more word for giving reasons.
\end{grammarlens}

\subsection*{Your turn}
\begin{itemize}
\raggedright
  \item \emph{Say it:} \emph{lìngwài}
  \item \emph{Say it:} \emph{lìngwài}, once more
  \item \emph{Say it:} \emph{zhège tài guì le lìngwài wǒ yě bù xǐhuan} — This one's too expensive, and besides, I don't like it.
  \item \emph{Recall:} say the Mandarin for to give an example, then the Mandarin for secondly, then say \emph{lìngwài} again
\end{itemize}

\begin{tcolorbox}[breakable,colback=teal!4,colframe=teal!35!black,title={Before you move on}]
What does lìngwài mean? (\textquotedblleft{}Besides; in addition\textquotedblright{}.) Say it once more.
\end{tcolorbox}
